\title{\color{ladrillo}\bfseries El Israel del Segundo Templo\\[4pt]\large\color{cafe}\mdseries Un retorno real pero incompleto: un pueblo sin rey que guarda la Torá y espera la reunión de todas sus tribus}
\author{}\date{\color{sepia}Serie «Israel, el pueblo de Dios» · Estudio 6}
\begin{document}
\maketitle\thispagestyle{empty}

\begin{abstract}\noindent
Este sexto estudio recorre la historia de Israel desde el edicto de Ciro hasta la ampliación del templo por Herodes: un retorno real pero incompleto, sin rey davídico, bajo imperios extranjeros y con la mayor parte del pueblo en la diáspora. Abarca de 538 a c.~20/19 a.\,e.\,c. Defiende dos tesis. La primera: en esta etapa Israel aprendió a ser «nación santa» sin trono y sin independencia, alrededor de la Torá leída en público, del sábado y de la sinagoga, y por medio de la Septuaginta puso la palabra de Dios al alcance de las naciones, aunque también conoció la tentación del nacionalismo y de la separación hostil. La segunda: el retorno de 538 fue una restauración real pero parcial, y los mismos textos de la época lo reconocen («he aquí que hoy somos siervos», Neh 9:36); por eso la esperanza profética de Isaías, Jeremías, Zacarías y Malaquías siguió abierta y se concentró en la reunión de todas las tribus y en la venida del Señor a su templo. El estudio distingue en cada caso entre «Israel» como pueblo completo, «Israel» como antiguo reino del norte y «Judá»; integra la arqueología persa y helenística (Cilindro de Ciro, tablillas de Al-Yahudu, papiros de Elefantina, sellos y monedas de Yehud, Garizim, Qumrán, inscripciones del templo herodiano), la tradición judía (Gran Asamblea, cese de la profecía, Amidá, Purim, Janucá, Seder Olam, Josefo) y corrige con respeto algunos datos del material de estudio «Octava lección». Concluye que Simeón y Ana, que esperaban «la consolación de Israel» y «la redención de Jerusalén» (Lc 2:25, 38), resumen la condición del pueblo al final del período.
\par\smallskip\noindent\textbf{Palabras clave:} Segundo Templo, retorno, Ciro, Esdras, Nehemías, diáspora, Septuaginta, macabeos, judíos, samaritanos, restauración, Romanos 11.
\end{abstract}

\section{Introducción}

¿Volvió Israel del exilio? La respuesta obvia es que sí: en 538 a.\,e.\,c. Ciro autorizó el regreso, el templo se reconstruyó y Jerusalén volvió a tener muros. Pero los libros de la época responden con más matices. Los ancianos que habían visto el primer templo lloraron al ver los cimientos del segundo (Esd 3:12); Nehemías, en plena restauración, ora: «he aquí que hoy somos siervos; en la tierra que diste a nuestros padres [\ldots] somos siervos» (Neh 9:36); y Daniel, al leer la profecía de los setenta años de Jeremías, recibe una respuesta que habla de «setenta semanas» (Dn 9:2, 24). Cinco siglos después, Simeón todavía esperaba «la consolación de Israel» (Lc 2:25), y los discípulos preguntaban a Jesús resucitado: «¿restaurarás el reino a Israel en este tiempo?» (Hch 1:6).

La tesis de este estudio es doble. Primero, el período del Segundo Templo es la escuela en que Israel aprendió a vivir como pueblo santo sin rey, sin independencia y en gran parte fuera de la tierra: la Torá leída y explicada en público, el sábado, la sinagoga, la oración diaria y la traducción de las Escrituras al griego sostuvieron su identidad y la abrieron a las naciones. Segundo, la restauración que comenzó con Ciro fue real pero incompleta. Volvieron sobre todo los de Judá y Benjamín, con sacerdotes, levitas y una parte de las otras tribus; la casa de David no recuperó el trono; la gloria no llenó el templo como en tiempos de Salomón; y los profetas de la época (Hageo, Zacarías, Malaquías) proyectaron la plena restauración hacia adelante. Esa esperanza abierta es la que recogen el Nuevo Testamento y, al final de la serie, Ro 11.

\begin{center}\small
\renewcommand{\arraystretch}{1.3}
\begin{tabular}{>{\RaggedRight\arraybackslash}p{3.1cm}>{\RaggedRight\arraybackslash}p{5.6cm}>{\RaggedRight\arraybackslash}p{5.2cm}}
\toprule
\textbf{Eje de la serie} & \textbf{Pregunta en esta etapa} & \textbf{Textos centrales}\\
\midrule
Nación santa y reino de sacerdotes & ¿Cómo vive su vocación un pueblo sin rey, sometido a imperios y disperso entre las naciones? & Esd 7:10; Neh 8:1--8; Zac 8:20--23; Is 56:6--8; Mal 1:11; 3:16--17\\
Promesa de restauración & ¿Qué parte de las promesas se cumplió en el retorno y qué quedó pendiente? & Esd 1:1--4; 6:17; Hag 2:6--9; Neh 9:36--37; Dn 9:24--27; Mal 3:1; 4:5--6 [3:23--24]\\
Proyección profética y paulina & ¿Cómo releen Isaías, Jeremías y el Nuevo Testamento este retorno incompleto? & Is 49:5--6; 60--62; 66:18--23; Jer 29:10--14; 33:14--18; Lc 2:25, 38; Hch 1:6; 26:7; Gá 4:4\\
\bottomrule
\end{tabular}\end{center}

\begin{nota}[Tres sentidos de «Israel» en este estudio]
Desde el Estudio 4 la serie distingue tres usos. (1) \textbf{Israel = todo el pueblo}, las doce tribus descendientes de Jacob. (2) \textbf{Israel = reino del norte} (diez tribus, Efraín, capital Samaria), deportado por Asiria en 722 a.\,e.\,c. (3) \textbf{Judá = reino del sur} (Judá y Benjamín, Jerusalén, casa de David), desterrado a Babilonia en 586 a.\,e.\,c. En el período persa aparece un cuarto término: \textbf{judíos} (\heb{יְהוּדִים}, \emph{yehudim}), que primero designa a los de Judá y luego a todo el pueblo que se identifica con el templo de Jerusalén, aunque sea de otra tribu (véase §4.2). Los libros de Esdras, Nehemías y Crónicas usan con frecuencia «Israel» en el sentido (1) para la comunidad que volvió: lo indicaremos en cada caso.
\end{nota}

\begin{nota}[Sobre traducción, numeración y fechas]
Las citas bíblicas son traducción propia desde el texto hebreo y arameo (BHS) y griego (NA28); la Septuaginta y los libros de Macabeos, Tobías y Ben Sirá se citan según Rahlfs-Hanhart. Cuando la numeración hebrea difiere de la española se da la española y la hebrea entre corchetes: Zac 2:10--12 [2:14--16], Neh 9:38 [10:1], Mal 4:5--6 [3:23--24]. Las fechas se indican como a.\,e.\,c. y e.\,c.
\end{nota}

\section{Marco cronológico}

El período tiene tres fases políticas. La \textbf{persa} va desde la toma de Babilonia por Ciro (539 a.\,e.\,c.) hasta la llegada de Alejandro Magno (332 a.\,e.\,c.); en ella se sitúan el retorno, la reconstrucción del templo (520--515), Ester, Esdras y Nehemías, y los últimos profetas canónicos. La \textbf{helenística} comienza con Alejandro y comprende el dominio de los Ptolomeos de Egipto (301--200 a.\,e.\,c.), el de los Seléucidas de Siria (desde 200 a.\,e.\,c.), la crisis de Antíoco IV y la rebelión macabea (167--164 a.\,e.\,c.) y el Estado asmoneo independiente (desde 142/140 a.\,e.\,c.). La \textbf{romana} empieza cuando Pompeyo toma Jerusalén (63 a.\,e.\,c.) e incluye el reinado de Herodes (37--4 a.\,e.\,c.), que amplió el templo desde c. 20/19 a.\,e.\,c. El Estudio~7 continúa desde el nacimiento de Jesús.

\begin{center}
\begin{tikzpicture}[x=1cm,y=1cm]
\draw[line width=1.6pt,cafe,-{Latex[length=3mm]}] (0,0) -- (15.4,0);
\foreach \x/\t/\b in {0.5/{Edicto de Ciro}/{538; Esd 1:1--4},
                      2.4/{Templo dedicado}/{515; Esd 6:15--17},
                      4.3/{Esdras y Nehemías}/{458 y 445; Neh 8},
                      6.2/{Alejandro}/{332},
                      8.1/{Seléucidas}/{200, Panias},
                      10.0/{Profanación y Janucá}/{167--164; 1~Mac 4:52},
                      11.9/{Asmoneos}/{142/140; 1~Mac 13:41},
                      13.6/{Pompeyo}/{63},
                      15.1/{Templo de Herodes}/{c. 20/19}}{
  \fill[oro] (\x,0) circle (4pt);
  \node[above,font=\scriptsize\bfseries,text=ladrillo,text width=1.7cm,align=center] at (\x,0.2) {\t};
  \node[below,font=\tiny,text=sepia,text width=1.7cm,align=center] at (\x,-0.2) {\b};}
\draw[decorate,decoration={brace,amplitude=5pt,mirror},sepia] (0.2,-0.85) -- (5.2,-0.85) node[midway,below=5pt,font=\scriptsize\itshape]{Período persa};
\draw[decorate,decoration={brace,amplitude=5pt,mirror},sepia] (5.7,-0.85) -- (12.7,-0.85) node[midway,below=5pt,font=\scriptsize\itshape]{Período helenístico y asmoneo};
\draw[decorate,decoration={brace,amplitude=5pt,mirror},sepia] (13.1,-0.85) -- (15.4,-0.85) node[midway,below=5pt,font=\scriptsize\itshape]{Roma};
\end{tikzpicture}\\[2pt]
{\footnotesize\color{sepia}\textit{Nota.} Elaboración propia. Fechas a.\,e.\,c.; varias están en debate (véase la nota siguiente).}
\end{center}

\begin{nota}[Fechas en debate]
\textbf{Dedicación del templo}: Esd 6:15 da el 3 de adar del sexto año de Darío~I, que la mayoría convierte en marzo de 515; algunos autores dan 516. \textbf{Esdras}: llegó «en el séptimo año del rey Artajerjes» (Esd 7:7); si es Artajerjes~I, en 458; si es Artajerjes~II, en 398, lo que lo pondría después de Nehemías. La Biblia los presenta juntos (Neh 8:9; 12:26), y este estudio sigue la fecha de 458, con Williamson (1985), señalando que el punto está abierto. \textbf{Nehemías}: el año veinte de Artajerjes~I (Neh 2:1), 445/444. \textbf{Macabeos}: la literatura antigua daba 168--165 para la profanación y la purificación; hoy se prefiere 167--164, a partir de la conversión de la era seléucida de 1~Mac y de fuentes cuneiformes (Bickerman, 1988; Goldstein, 1976). \textbf{Herodes}: Josefo da el año 18 del reinado (\emph{Ant.} 15.380, c. 20/19) y, en otro lugar, el año 15 (\emph{Guerra} 1.401).
\end{nota}

\section{Contexto histórico y arqueológico}

\subsection{Babilonia y Persia: el Cilindro de Ciro y Al-Yahudu}
El \textbf{Cilindro de Ciro}, hallado en Babilonia en 1879 y conservado en el Museo Británico, es una proclamación en acadio redactada tras la toma de la ciudad en 539 a.\,e.\,c. Ciro se presenta como elegido por Marduk y dice haber devuelto a sus santuarios las imágenes de dioses de varias ciudades «al otro lado del Tigris» y haber reunido a sus habitantes para hacerlos volver a sus moradas (Finkel, 2013). El texto ilustra la política que Esd 1:1--4 atribuye a Ciro, pero \destacar{no menciona a Judá, a Jerusalén ni a YHWH}: las ciudades que nombra están en Mesopotamia y el oeste de Irán. No es, por tanto, el edicto de Esdras, sino un testimonio paralelo de la misma política (Kuhrt, 1983). Esd 6:2--5 conserva además un memorándum arameo hallado en Ecbatana con medidas del templo y la devolución de los utensilios.

En Babilonia, más de cien tablillas cuneiformes publicadas procedentes de \textbf{Al-Yahudu}, «la ciudad de Judá» o «de los judeos», y de aldeas vecinas documentan la vida de los deportados entre 572 y 477 a.\,e.\,c.: contratos, arriendos, deudas, con nombres teóforos como \emph{Yahu-azar} o \emph{Ahi-Yahu} (Pearce y Wunsch, 2014). Muestran que los desterrados de Judá conservaron su identidad y que muchos se instalaron de forma estable: la carta de Jer 29:5--7 («edificad casas [\ldots] procurad la paz de la ciudad») describe exactamente esa situación. Explican también por qué la mayoría no volvió con Ciro.

\subsection{Egipto: Elefantina y Alejandría}
En la isla de \textbf{Elefantina}, en la primera catarata del Nilo, una guarnición judía al servicio de Persia tenía desde antes de Cambises (525 a.\,e.\,c.) un templo a YHW. Sus papiros arameos llaman a la comunidad «la guarnición judía». El más conocido es la \textbf{carta a Bagohi}, gobernador de Yehud, fechada en 407 a.\,e.\,c.: los sacerdotes de Elefantina cuentan que los sacerdotes del dios Jnum destruyeron su templo en 410, piden ayuda para reconstruirlo, dicen haber escrito antes sin respuesta al sumo sacerdote Johanán de Jerusalén (cf. Neh 12:22) y mencionan a «Delaías y Selemías, hijos de Sanbalat, gobernador de Samaria» (Porten y Yardeni, 1986). La carta confirma la existencia histórica del Sanbalat que se opuso a Nehemías (Neh 2:10; 4:1--2 [3:33--34]). Otro papiro, de 419/418, transmite instrucciones sobre la fiesta de los panes sin levadura. Elefantina muestra una diáspora fiel al Dios de Israel pero con prácticas que el judaísmo de Jerusalén no aprobaba (un templo fuera de Jerusalén, y menciones de otras divinidades junto a YHW).

\textbf{Alejandría}, fundada por Alejandro en 331 a.\,e.\,c., llegó a tener una de las mayores comunidades judías del mundo antiguo. Allí se tradujo la Torá al griego en el siglo \textsc{iii} a.\,e.\,c. y allí escribiría Filón en el siglo \textsc{i} e.\,c. En Leontópolis, Onías~IV, hijo del sumo sacerdote depuesto, fundó hacia 160 a.\,e.\,c. otro templo que funcionó hasta 73 e.\,c. (Josefo, \emph{Ant.} 13.62--73).

\subsection{Yehud y Samaria}
La provincia persa de Judá se llamaba \emph{Yehud} (arameo \heb{יְהוּד}; Esd 5:8, «la provincia de Judea»). Cientos de \textbf{impresiones de sello} con las letras \emph{yhd} o \emph{yhwd} en asas de jarras, y pequeñas \textbf{monedas de plata} del siglo \textsc{iv} con la leyenda \emph{yhd}, algunas con el nombre de «Ezequías el gobernador» (\emph{Yeḥizqiyyah ha-peḥah}) o de «Johanán el sacerdote», muestran una provincia pequeña con administración propia (Lipschits y Vanderhooft, 2011; Meshorer, 2001). La población de Jerusalén era reducida, lo que concuerda con Neh 7:4 («la ciudad era espaciosa y grande, pero poco pueblo dentro de ella») y con el sorteo para repoblarla (Neh 11:1).

\begin{nota}[Debate arqueológico: el muro de Nehemías]
Eilat Mazar, en sus excavaciones de 2005--2007 en la Ciudad de David, atribuyó a Nehemías una torre y un tramo de muro sobre la ladera oriental. Finkelstein (2008) sostiene que esas estructuras son asmoneas (siglo \textsc{ii} a.\,e.\,c.) y que la Jerusalén persa era tan pequeña que el relato de Neh 3 refleja, en parte, una realidad posterior. Otros responden que Nehemías reparó sobre todo trazados antiguos. Es un debate abierto: no hay hoy un tramo de muro que pueda atribuirse con seguridad a Nehemías.
\end{nota}

En el monte \textbf{Garizim}, junto a Siquem, las excavaciones de Yitzhak Magen revelaron un recinto sagrado samaritano construido a mediados del siglo \textsc{v} a.\,e.\,c., en época persa, ampliado en época helenística y destruido por Juan Hircano hacia 112/111 a.\,e.\,c. (Magen, 2008). Josefo situaba su fundación en tiempos de Alejandro (\emph{Ant.} 11.302--347); la arqueología la adelanta más de un siglo. El hallazgo da contexto a Jn 4:20: «nuestros padres adoraron en este monte».

\subsection{Del helenismo a Herodes}
Del período helenístico y romano proceden varios hallazgos directamente relacionados con la historia bíblica:
\begin{itemize}
\item \textbf{Estela de Heliodoro} (hallada en Maresá, publicada en 2007): carta de Seleuco~IV a su ministro Heliodoro (178 a.\,e.\,c.) sobre el control de los santuarios de Celesiria y Fenicia (Cotton y Wörrle, 2007). Da contexto al episodio de 2~Mac 3 y a Dn 11:20 («uno que hará pasar un cobrador de tributos»).
\item \textbf{Rollos del Mar Muerto} (Qumrán, desde 1947): unos 900 manuscritos, entre ellos el gran rollo de Isaías \textbf{1QIsa\textsuperscript{a}}, copiado hacia el siglo \textsc{ii} a.\,e.\,c., con el libro completo y diferencias respecto al Texto Masorético que en su mayoría son ortográficas (Ulrich y Flint, 2010). Los rollos documentan la diversidad del judaísmo y su esperanza mesiánica (§7.4).
\item \textbf{Muro herodiano}: los sillares con márgenes rebajados del monte del templo. En la esquina suroeste, Benjamín Mazar halló en 1968--1969 una piedra caída con la inscripción hebrea \emph{le-beit ha-teqi'ah le-hak[riz]}, «al lugar de tocar la trompeta para anun[ciar]», que corresponde al sacerdote que anunciaba con trompeta el comienzo y el fin del sábado (Josefo, \emph{Guerra} 4.582; Mazar, 1975).
\item \textbf{Inscripción de advertencia del templo}: Charles Clermont-Ganneau halló en 1871 una placa griega completa (un fragmento de otra apareció en 1935) que prohibía a los extranjeros pasar la balaustrada que rodeaba el santuario: «Ningún extranjero entre dentro de la balaustrada y del recinto que rodea el santuario; quien sea sorprendido será responsable de su propia muerte». Josefo describe estas placas en griego y en latín (\emph{Guerra} 5.193--194). Es el trasfondo de Hch 21:28--29 y del «muro intermedio de separación» de Ef 2:14.
\item \textbf{Inscripción de Teodoto} (Ofel, 1913--1914): Teodoto, sacerdote y jefe de sinagoga, hijo y nieto de jefes de sinagoga, construyó en Jerusalén una sinagoga «para la lectura de la ley y la enseñanza de los mandamientos», con hospedaje para los que venían del extranjero. Se fecha en general antes de 70 e.\,c. (Kloppenborg, 2000) y recuerda la sinagoga de los libertos de Hch 6:9.
\end{itemize}

\section{Desarrollo bíblico}

\subsection{El edicto de Ciro y el primer retorno}
La Biblia narra el retorno como cumplimiento de profecía. Crónicas termina y Esdras comienza con el mismo edicto:
\begin{cita}En el primer año de Ciro, rey de Persia, para que se cumpliera la palabra de YHWH por boca de Jeremías, YHWH despertó el espíritu de Ciro [\ldots]: «YHWH, el Dios de los cielos, me ha dado todos los reinos de la tierra, y me ha encargado que le edifique una casa en Jerusalén, que está en Judá. Quien haya entre vosotros de todo su pueblo, sea su Dios con él, y suba».\fuente{(Esd 1:1--3; 2 Cr 36:22--23)}\end{cita}

La «palabra de Jeremías» son los setenta años (Jer 25:11--12; 29:10). Isaías había nombrado a Ciro como instrumento de Dios:
\hebverso{כֹּה־אָמַר יְהוָה לִמְשִׁיחוֹ לְכוֹרֶשׁ אֲשֶׁר־הֶחֱזַקְתִּי בִימִינוֹ}
\begin{cita}Así dice YHWH a su ungido, a Ciro, a quien tomé de la mano derecha.\fuente{(Is 45:1; cf. 44:28)}\end{cita}
Es el único lugar de la Biblia hebrea en que un rey extranjero recibe el título de \heb{מָשִׁיחַ} (\emph{mashíaj}, «ungido»). El título describe una función (un rey puesto por Dios para liberar a su pueblo), no un mesías davídico. Para la mayoría de la investigación, Is 40--55 procede del final del exilio; la tradición judía y cristiana lo atribuye al profeta del siglo \textsc{viii}. En cualquiera de las dos lecturas, el texto interpreta a Ciro como siervo del plan de Dios.

El primer grupo sube bajo \textbf{Sesbasar}, «príncipe de Judá» (Esd 1:8), que lleva los utensilios del templo y pone los cimientos (Esd 5:16). Poco después aparecen como líderes \textbf{Zorobabel}, nieto del rey Joaquín (1~Cr 3:17--19, con una dificultad genealógica sobre su padre), y el sumo sacerdote \textbf{Josué} (Jesúa). Levantan el altar y celebran Sucot (Esd 3:1--6). Al poner los cimientos, los jóvenes gritan de alegría y los ancianos lloran (Esd 3:12--13). La oposición de los vecinos detiene la obra (Esd 4) hasta el segundo año de Darío~I (520 a.\,e.\,c.).

\subsection{¿Quiénes vuelven? Israel, Judá y los «judíos»}
El texto bíblico es preciso. Los que responden al edicto son «los jefes de las casas paternas de Judá y de Benjamín, los sacerdotes y los levitas» (Esd 1:5): es decir, sobre todo el antiguo \textbf{reino del sur}. Pero la lista de Esd 2 se titula «el número de los varones del pueblo de \emph{Israel}» (Esd 2:2) y termina diciendo que «todo Israel» habitó en sus ciudades (Esd 2:70): aquí «Israel» tiene el sentido (1), el pueblo entero, representado por los que volvieron. El mismo uso aparece en la dedicación del templo:

\hebverso{וצפירי עזין לחטאה על כל ישראל תרי עשר למנין שבטי ישראל}
\begin{cita}Y como ofrenda por el pecado por todo Israel, doce machos cabríos, conforme al número de las tribus de Israel.\fuente{(Esd 6:17, arameo)}\end{cita}

La comunidad que volvió, mayoritariamente de Judá y Benjamín, ofrece por las \emph{doce} tribus. No pretende ser todo Israel, sino actuar en nombre de todo Israel. Esdras repite el gesto con «doce becerros por todo Israel» y doce machos cabríos (Esd 8:35). Otros textos muestran que en la comunidad había también miembros de las tribus del norte: «en Jerusalén habitaron de los hijos de Judá, de los hijos de Benjamín, y de los hijos de Efraín y de Manasés» (1~Cr 9:3), y Neh 11:20 habla del «resto de Israel» en las ciudades de Judá. Esto concuerda con lo visto en el Estudio~5: antes del exilio, grupos del norte se habían unido a Judá bajo Asa, Ezequías y Josías (2~Cr 15:9; 30:11; 34:9). Siglos después, Lucas presenta a Ana, profetisa, «de la tribu de Aser» (Lc 2:36), una tribu del norte, y Pablo dice ser «de la tribu de Benjamín» (Ro 11:1; Flp 3:5).

\begin{nota}[Precisión: el alcance de 1 Cr 9:3]
1~Cr 9:2--34 tiene un paralelo parcial en Neh 11:3--19, pero Neh 11 menciona solo a Judá y Benjamín. Algunos comentaristas leen la mención de Efraín y Manasés en 1~Cr 9:3 como reflejo de la presencia real de norteños en Jerusalén; otros, como la visión del Cronista, que piensa siempre en «todo Israel». En cualquier caso, la Biblia no dice que las diez tribus volvieran como tales: dice que el núcleo del retorno fue Judá y Benjamín y que hubo norteños integrados en él.
\end{nota}

\textbf{El nombre «judíos».} \heb{יְהוּדִי} (\emph{yehudí}) significa en origen «de Judá». Aparece ya antes del exilio (2~R 16:6; Jer 34:9), pero se generaliza en el período persa, cuando la provincia se llama \emph{Yehud} y sus habitantes \emph{yehudaye} en arameo (Esd 4:12; 5:1). El término deja entonces de ser solo tribal: Mardoqueo es «un judío (\heb{אִישׁ יְהוּדִי}) [\ldots] benjaminita (\heb{אִישׁ יְמִינִי})» (Est 2:5), y Zacarías anuncia que las naciones se asirán del manto de «un judío» (Zac 8:23). Josefo lo explica así: se llaman judíos «desde el día en que subieron de Babilonia, por la tribu de Judá, que llegó primero a estos lugares» (\emph{Ant.} 11.173). En el Nuevo Testamento \grc{Ἰουδαῖος} designa a todo el pueblo, y Pablo puede decir a la vez «soy judío» (Hch 21:39; 22:3) y «soy israelita, de la tribu de Benjamín» (Ro 11:1).

\textbf{Los samaritanos.} Cuando los vecinos del norte se ofrecen a construir con los repatriados («buscamos a vuestro Dios como vosotros, y a él sacrificamos desde los días de Esar-hadón, rey de Asiria, que nos hizo subir aquí», Esd 4:2), Zorobabel los rechaza (Esd 4:3). El texto los llama «enemigos de Judá y de Benjamín» (Esd 4:1) y los vincula con los colonos traídos por Asiria (2~R 17:24--41). Los samaritanos, por su parte, se consideran descendientes de Efraín y Manasés y guardianes (\emph{shamerim}) de la Torá, y conservan su propio Pentateuco. La realidad probable es mixta: población israelita del antiguo reino del norte que no fue deportada, unida a colonos extranjeros (Crown, 1989). El conflicto se agravó con Nehemías y Sanbalat y culminó con el templo del Garizim. Es una herida dentro de la casa de Jacob que el Nuevo Testamento vuelve a abrir para sanarla (Lc 10:33; Jn 4; Hch 1:8; 8:5--25).

\subsection{Hageo: la gloria postrera de esta casa}
En 520 a.\,e.\,c. Hageo reprocha al pueblo vivir en «casas artesonadas» mientras la casa de Dios está en ruinas (Hag 1:4). La obra se reanuda. Ante la decepción de los que recordaban el primer templo (Hag 2:3), el profeta promete:

\hebverso{גָּדוֹל יִהְיֶה כְּבוֹד הַבַּיִת הַזֶּה הָאַחֲרוֹן מִן־הָרִאשׁוֹן}
\begin{cita}Haré temblar los cielos y la tierra [\ldots] y haré temblar a todas las naciones, y vendrá lo precioso de todas las naciones, y llenaré de gloria esta casa [\ldots]. Mayor será la gloria de esta casa postrera que la de la primera, dice YHWH de los ejércitos; y en este lugar daré paz.\fuente{(Hag 2:6--9)}\end{cita}

\begin{nota}[Precisión filológica: \emph{ḥemdat kol-ha-goyim} (Hag 2:7)]
\heb{חֶמְדַּת כָּל־הַגּוֹיִם} es literalmente «lo deseable de todas las naciones», con verbo en plural (\emph{u-va'u}, «vendrán»). La mayoría de los intérpretes modernos lo entiende como los tesoros de las naciones que llegan al templo (cf. Hag 2:8; Is 60:5--11). La Vulgata tradujo \emph{desideratus cunctis gentibus}, «el deseado de todas las naciones», y la lectura mesiánica se volvió tradicional en el cristianismo. Ambas lecturas sitúan la gloria futura del templo en relación con las naciones.
\end{nota}

Hageo termina con una palabra a Zorobabel: «te tomaré [\ldots] y te pondré como anillo de sellar» (Hag 2:23). Es la inversión exacta de Jer 22:24, donde Dios decía que arrancaría a Jeconías, abuelo de Zorobabel, aunque fuera su anillo. La línea davídica, cortada del trono, sigue elegida. Sin embargo, Zorobabel nunca fue rey y desaparece del relato sin explicación.

\subsection{Zacarías: Judá, Efraín y las naciones}
Zacarías comienza su ministerio dos meses después de Hageo (Zac 1:1). Su primera visión recoge la pregunta del período: «¿hasta cuándo no tendrás piedad de Jerusalén [\ldots] con las cuales has estado airado estos setenta años?» (Zac 1:12). Varios de sus oráculos son centrales para esta serie:
\begin{itemize}
\item \textbf{Zac 2:10--12 [2:14--16]}: «Canta y alégrate, hija de Sion, porque vengo y habitaré en medio de ti [\ldots]. Y se unirán muchas naciones a YHWH en aquel día, y me serán por pueblo [\ldots]. YHWH poseerá a Judá, su porción en la tierra santa, y escogerá aún a Jerusalén». Es la única vez que la Biblia hebrea dice «tierra santa» (\heb{אַדְמַת הַקֹּדֶשׁ}). Las naciones se unen a YHWH sin que Judá deje de ser su porción.
\item \textbf{Zac 8:13}: «como fuisteis maldición entre las naciones, \emph{oh casa de Judá y casa de Israel}, así os salvaré y seréis bendición». Después del retorno, el profeta sigue nombrando a las dos casas: la restauración incluye al reino del norte disperso. Lo mismo en Zac 9:13 («Judá [\ldots] Efraín»), 10:6 («la casa de Judá [\ldots] la casa de José») y 11:14, donde la «hermandad entre Judá e Israel» se rompe simbólicamente.
\item \textbf{Zac 8:20--23}: diez hombres de todas las lenguas se asirán del manto de un judío (texto en §5).
\item \textbf{Zac 9:9}: «He aquí tu rey viene a ti, justo y salvador, humilde (\heb{עָנִי}), y cabalgando sobre un asno». Un rey sin ejército para un pueblo sin rey (cf. Mt 21:5; Jn 12:15).
\item \textbf{Zac 12:10}: «derramaré sobre la casa de David y sobre los moradores de Jerusalén espíritu de gracia y de oración, y mirarán a mí, a quien traspasaron (\heb{וְהִבִּיטוּ אֵלַי אֵת אֲשֶׁר־דָּקָרוּ}), y llorarán por él». Juan lo aplica a la cruz (Jn 19:37); muchos intérpretes lo relacionan con el futuro de Israel que Pablo anuncia en Ro 11:26.
\end{itemize}
La crítica atribuye Zac 9--14 a una mano posterior; la tradición, al mismo profeta. Para este estudio basta notar que el libro, tal como fue recibido en el Segundo Templo, une restauración de las dos casas, inclusión de las naciones y un Mesías humilde y traspasado.

El templo se terminó el 3 de adar del sexto año de Darío (Esd 6:15). La comunidad celebró la Pascua: «los hijos de Israel que habían vuelto del cautiverio y todos los que se habían apartado de la inmundicia de las naciones de la tierra para buscar a YHWH, Dios de Israel» (Esd 6:21). La Pascua incluye a quienes se unen al pueblo, como en Éx 12:48.

\subsection{Ester y Purim: el pueblo esparcido}
El libro de Ester transcurre en Susa bajo Asuero, que suele identificarse con Jerjes~I (486--465 a.\,e.\,c.). Amán describe a Israel con precisión: «hay un pueblo esparcido y dividido entre los pueblos en todas las provincias de tu reino, y sus leyes son diferentes» (Est 3:8). La liberación da origen a Purim, el 14 y 15 de adar (Est 9:20--32). El libro tiene dos rasgos notables: el texto hebreo no menciona a Dios, y al final «muchos de los pueblos de la tierra se hacían judíos» (\heb{מִתְיַהֲדִים}, \emph{mityahadim}, Est 8:17), única aparición del verbo en la Biblia; la Septuaginta traduce «se circuncidaban y judaizaban». La historicidad de algunos detalles (la reina Ester frente a la Amestris de Heródoto) se discute; la descripción de la corte persa es, en cambio, muy precisa. 2~Mac 15:36 llama ya al 14 de adar «el día de Mardoqueo», la mención más antigua de Purim fuera de Ester.

\subsection{Esdras y Nehemías: Torá, muros y pacto}
Esdras es «escriba versado en la ley de Moisés» (Esd 7:6), y su programa se resume en un versículo:
\begin{cita}Porque Esdras había preparado su corazón para inquirir la ley de YHWH y para cumplirla, y para enseñar en Israel estatutos y decretos.\fuente{(Esd 7:10)}\end{cita}
El verbo es \heb{לִדְרוֹשׁ} (\emph{lidrosh}), de donde vendrá \emph{midrash}. El orden es estudiar, cumplir, enseñar. Nehemías, copero de Artajerjes~I, obtiene permiso para reconstruir los muros, que se terminan en 52 días (Neh 6:15) pese a la oposición de Sanbalat de Samaria, Tobías el amonita y Gesem el árabe (Neh 2:19).

El centro de ambos libros es la lectura pública de la Torá el primer día del mes séptimo (Neh 8):
\hebverso{וַיִּקְרְאוּ בַסֵּפֶר בְּתוֹרַת הָאֱלֹהִים מְפֹרָשׁ וְשׂוֹם שֶׂכֶל וַיָּבִינוּ בַּמִּקְרָא}
\begin{cita}Y leían en el libro de la ley de Dios con claridad, y ponían el sentido, de modo que entendían la lectura.\fuente{(Neh 8:8)}\end{cita}
Hombres, mujeres «y todos los que podían entender» (Neh 8:2) escuchan desde el alba hasta el mediodía; los levitas explican. El pueblo llora, y Nehemías y Esdras le dicen: «el gozo de YHWH es vuestra fuerza» (Neh 8:10). Después celebran Sucot como no se hacía «desde los días de Josué hijo de Nun» (Neh 8:17), confiesan los pecados de los padres (Neh 9) y firman un pacto (\heb{אֲמָנָה}, \emph{amaná}, Neh 9:38 [10:1]): no dar sus hijas a los pueblos vecinos, guardar el sábado y el año sabático, sostener el templo (Neh 10). El pueblo se reconstituye en torno a un libro.

Las reformas incluyeron la disolución de matrimonios mixtos (Esd 9--10; Neh 13:23--27). El texto motiva la medida por la idolatría (Esd 9:1; Neh 13:26, con el ejemplo de Salomón), no por la etnia; prosélitos como Rut, antepasada de David, muestran que la Torá admitía la incorporación de extranjeros que adoptaban al Dios de Israel. Aun así, esta línea de separación se endurecerá en siglos posteriores (§5).

\subsection{Malaquías: el mensajero y Elías}
Malaquías, el último de los Doce, refleja la situación posterior a la reconstrucción: sacrificios defectuosos, sacerdotes negligentes, divorcios, diezmos retenidos. Su mensaje se abre con el amor de Dios por Jacob (Mal 1:2) y contiene una afirmación sobre las naciones que sorprende en ese contexto: «desde donde nace el sol hasta donde se pone, es grande mi nombre entre las naciones, y en todo lugar se ofrece a mi nombre incienso y ofrenda limpia» (Mal 1:11). El libro anuncia la venida del Señor a su templo:
\hebverso{הִנְנִי שֹׁלֵחַ מַלְאָכִי וּפִנָּה־דֶרֶךְ לְפָנָי}
\begin{cita}He aquí que yo envío mi mensajero, y preparará el camino delante de mí; y vendrá súbitamente a su templo el Señor a quien vosotros buscáis, y el mensajero del pacto, a quien vosotros deseáis.\fuente{(Mal 3:1)}\end{cita}
Y funda la permanencia de Israel en la inmutabilidad de Dios:
\hebverso{כִּי אֲנִי יְהוָה לֹא שָׁנִיתִי וְאַתֶּם בְּנֵי־יַעֲקֹב לֹא כְלִיתֶם}
\begin{cita}Porque yo, YHWH, no cambio; por eso vosotros, hijos de Jacob, no habéis sido consumidos.\fuente{(Mal 3:6)}\end{cita}
\noindent El canon de los profetas se cierra con dos mandatos: «acordaos de la ley de Moisés» (Mal 4:4 [3:22]) y la promesa de Elías, que «hará volver el corazón de los padres hacia los hijos» antes del día de YHWH (Mal 4:5--6 [3:23--24]). Lucas aplica el texto a Juan el Bautista (Lc 1:17).

\subsection{Los «cuatrocientos años de silencio»}
La expresión es una convención cristiana para el tiempo entre Malaquías y Juan el Bautista. Es exacta en un sentido: no hay profeta canónico, y la propia literatura judía lo reconoce (§6.2). Pero no fue un tiempo vacío. Fue el período más agitado de la historia del Segundo Templo.

\textbf{Alejandro y los reinos helenísticos.} Alejandro venció a Persia en Iso (333) y tomó la costa en 332. Josefo cuenta que, al acercarse a Jerusalén, Alejandro se postró ante el sumo sacerdote Jadúa y que le mostraron el libro de Daniel (\emph{Ant.} 11.329--339); el Talmud tiene una versión con Simón el Justo (Yoma 69a). La historia es probablemente una leyenda: ninguna fuente griega la recoge. Tras la muerte de Alejandro (323), Judea quedó bajo los \textbf{Ptolomeos} de Egipto (301--200) y, después de la batalla de Panias (200), bajo los \textbf{Seléucidas} de Siria. Dn 11:5--20 describe con detalle las guerras entre «el rey del sur» y «el rey del norte».

\textbf{La Septuaginta.} La \emph{Carta de Aristeas} cuenta que Ptolomeo~II (285--246) pidió al sumo sacerdote setenta y dos sabios, seis por tribu, que tradujeron la Torá en setenta y dos días. El relato es legendario y apologético; el Talmud tiene una versión propia (Meguilá 9a--b), y Filón dice que en su tiempo se celebraba una fiesta anual en la isla de Faros por la traducción (\emph{Vida de Moisés} 2.41). Lo histórico es que el Pentateuco se tradujo al griego en Alejandría en el siglo \textsc{iii} a.\,e.\,c., y que hacia 132 a.\,e.\,c. el nieto de Ben Sirá conocía ya en griego «la ley, los profetas y los demás libros» (prólogo de Eclesiástico) (Jobes y Silva, 2015). El nombre de «setenta» (LXX) procede de la leyenda.

\textbf{La helenización.} Bajo Antíoco~IV Epífanes (175--164), hijo de Antíoco~III el Grande, Jasón compró el sumo sacerdocio, construyó un gimnasio en Jerusalén y promovió costumbres griegas (2~Mac 4:7--15). 2~Mac 4:13 lo llama «auge del helenismo» (\grc{ἀκμή τις Ἑλληνισμοῦ}); es en este contexto donde aparece por primera vez la palabra \grc{Ἰουδαϊσμός}, «judaísmo» (2~Mac 2:21; 8:1), como lo opuesto. Menelao desplazó a Jasón ofreciendo más dinero, y en 171/170 a.\,e.\,c. hizo asesinar al legítimo sumo sacerdote \textbf{Onías~III} en Dafne, junto a Antioquía, por medio de Andrónico; Antíoco, al volver, mandó ejecutar a Andrónico (2~Mac 4:33--38).

\textbf{Antíoco~IV y la abominación.} En 169 Antíoco saqueó el templo (1~Mac 1:20--24). En 168 los romanos lo obligaron a retirarse de Egipto (el episodio del círculo trazado en la arena por Popilio Lenas en Eleusis, cerca de Alejandría; cf. Dn 11:30, «naves de Quitim»). En 167 prohibió la circuncisión, el sábado y la Torá, y el 15 de kislev del año 145 seléucida (diciembre de 167) se erigió sobre el altar «la abominación de la desolación» (\grc{βδέλυγμα ἐρημώσεως}, 1~Mac 1:54); el 25 de kislev se sacrificó por primera vez sobre el altar pagano (1~Mac 1:59). 2~Mac 6:2 dice que el templo fue dedicado a «Zeus Olímpico» (\grc{Διὸς Ὀλυμπίου}) y el del Garizim a «Zeus Hospitalario». Daniel había descrito al «cuerno pequeño» que quita el sacrificio continuo (Dn 8:9--12) y pone «la abominación desoladora»:
\hebverso{וְהֵסִירוּ הַתָּמִיד וְנָתְנוּ הַשִּׁקּוּץ מְשׁוֹמֵם}
\begin{cita}Y quitarán el continuo sacrificio, y pondrán la abominación desoladora.\fuente{(Dn 11:31; cf. 9:27; 12:11)}\end{cita}

\begin{nota}[Precisión: la fecha del libro de Daniel]
La tradición judía y cristiana sitúa a Daniel en el siglo \textsc{vi} a.\,e.\,c.; la mayoría de la investigación crítica fecha la redacción final de Dn 7--12 hacia 167--164 a.\,e.\,c., por la exactitud con que describe la crisis de Antíoco (Collins, 1993). Este estudio no resuelve el debate: lo relevante es que los judíos de la época leyeron su crisis a la luz de Daniel, y que Jesús retomó la expresión «abominación de la desolación» para un acontecimiento futuro (Mt 24:15; Mc 13:14).
\end{nota}

\textbf{Los macabeos y Janucá.} En Modín, el sacerdote Matatías se negó a sacrificar y huyó a los montes con sus cinco hijos (1~Mac 2). Su hijo Judas, llamado Macabeo, venció a los generales seléucidas y recuperó Jerusalén. El 25 de kislev del año 148 seléucida (diciembre de 164 a.\,e.\,c.), tres años después de la profanación, el templo fue purificado, se levantó un altar nuevo y se celebró su dedicación durante ocho días (1~Mac 4:36--59; 2~Mac 10:1--8). Esa fiesta de la dedicación (\heb{חֲנֻכָּה}, \emph{janucá}) es «la fiesta de la dedicación» en que Jesús camina por el pórtico de Salomón (\grc{τὰ ἐγκαίνια}, Jn 10:22).

\begin{nota}[Precisión respetuosa sobre la «Octava lección»]
El material de estudio «Israel, el pueblo de Dios» (Octava lección) presenta bien el sentido de esta crisis, pero contiene algunos datos que conviene ajustar:
\begin{itemize}
\item \textbf{Filiación.} Antíoco~IV era hijo de \emph{Antíoco~III} el Grande y hermano de Seleuco~IV, no hijo de Filipo de Macedonia (padre de Alejandro, muerto en 336 a.\,e.\,c.).
\item \textbf{Orden de los dominios.} Tras Alejandro, Judea estuvo primero bajo los Ptolomeos de Egipto (301--200) y después bajo los Seléucidas de Siria, no al revés. Ptolomeo~VI Filométor y su rival Ptolomeo~VIII eran hermanos (no padre e hijos), ambos sobrinos de Antíoco.
\item \textbf{La fecha de la reinauguración.} Fue el 25 de kislev de 164 a.\,e.\,c. (1~Mac 4:52), no de 139; la cifra 165 de obras antiguas responde a otra conversión de la era seléucida.
\item \textbf{«Júpiter».} 2~Mac 6:2 dice Zeus Olímpico; «Júpiter» es su equivalente latino (Vulgata). 1~Mac 1:54, 59 habla de una «abominación» y un altar pagano \emph{sobre el altar} de los holocaustos; no dice que se pusiera una imagen en el Lugar Santísimo. El sacrificio de cerdos sobre el altar lo menciona Josefo (\emph{Ant.} 12.253).
\item \textbf{Onías~III.} Fue asesinado c. 171/170 a.\,e.\,c. por Andrónico a instancias de Menelao (2~Mac 4:33--35), no por mano directa de Antíoco. Muchos intérpretes lo relacionan con Dn 9:26 («se quitará la vida al ungido») y Dn 11:22 («el príncipe del pacto»); el «príncipe del ejército» de Dn 8:11 suele entenderse como Dios mismo.
\item \textbf{Las 2.300 «tardes y mañanas».} La lección acierta al señalar que el hebreo de Dn 8:14 no dice «días». Su equivalencia con 1.150 días es una lectura posible (contar por separado el sacrificio de la tarde y el de la mañana), pero no la única: otros la leen como 2.300 días completos (unos seis años y cuatro meses, desde el comienzo de la crisis c. 171) y la interpretación historicista como 2.300 años. Los datos de 1~Mac van del 15 de kislev de 167 al 25 de kislev de 164: tres años y diez días, algo menos de 1.150 días (Collins, 1993).
\end{itemize}
\end{nota}

\textbf{Los asmoneos.} Los hermanos de Judas continuaron la lucha. Jonatán fue nombrado sumo sacerdote (152) y Simón obtuvo la independencia en 142: «fue quitado el yugo de las naciones de Israel» (1~Mac 13:41). En 140 una asamblea lo proclamó jefe y sumo sacerdote «para siempre, hasta que se levante un profeta fiel» (1~Mac 14:41). Juan Hircano~I (134--104) conquistó Samaria, destruyó el templo del Garizim y obligó a los idumeos a circuncidarse (Josefo, \emph{Ant.} 13.254--258). Aristóbulo~I tomó el título de rey; Alejandro Janeo (103--76) amplió el territorio y crucificó a ochocientos opositores (\emph{Ant.} 13.380). Una dinastía sacerdotal no davídica gobernaba como rey en Jerusalén, y no todos lo aceptaron (Schürer, 1973--1987).

\textbf{Fariseos, saduceos y esenios.} Josefo menciona las tres «escuelas» por primera vez en tiempos de Jonatán (\emph{Ant.} 13.171--173). Los \textbf{fariseos} (de \emph{perushim}, «separados») enseñaban la tradición de los padres además de la Torá escrita y creían en la resurrección; los \textbf{saduceos}, vinculados a la aristocracia sacerdotal, aceptaban solo la ley escrita y negaban la resurrección (Hch 23:8); los \textbf{esenios} vivían en comunidades de pureza estricta (\emph{Guerra} 2.119--161). Muchos investigadores identifican con ellos a la comunidad de Qumrán, aunque la identificación se discute (VanderKam, 2010).

\textbf{Roma y Herodes.} La guerra civil entre Hircano~II y Aristóbulo~II abrió la puerta a Pompeyo, que tomó Jerusalén en 63 a.\,e.\,c. y entró en el Lugar Santísimo, sin tocar sus tesoros (\emph{Ant.} 14.71--72). Herodes, idumeo de origen, fue nombrado rey por el Senado romano en 40 y tomó Jerusalén en 37. Hacia 20/19 a.\,e.\,c. emprendió la ampliación del templo y de la explanada, una obra que seguía en curso en tiempos de Jesús: «en cuarenta y seis años fue edificado este templo» (Jn 2:20). Hageo había prometido que la gloria de esa casa sería mayor que la del primer templo; en tamaño y belleza, el templo de Herodes lo superó, pero la plenitud de la promesa quedaba aún por venir.

\section{Israel como nación santa y reino de sacerdotes}

La vocación de Éx 19:5--6 se vivió en esta etapa en condiciones nuevas. Israel ya no tenía rey davídico ni, la mayor parte del tiempo, independencia; la mayoría de sus miembros vivía fuera de la tierra. Lo que lo mantuvo como «nación santa» fueron otros medios:

\begin{itemize}
\item \textbf{La Torá en el centro.} Desde Neh 8, el pueblo se define por un libro leído, explicado y obedecido. Esdras «enseña en Israel» (Esd 7:10) y los levitas «hacen entender» (Neh 8:7). Malaquías recuerda que «los labios del sacerdote han de guardar la sabiduría» y que el sacerdote es «mensajero de YHWH» (Mal 2:7).
\item \textbf{El sábado como frontera.} Nehemías cierra las puertas de Jerusalén el sábado (Neh 13:15--22); bajo Antíoco, muchos prefieren morir antes que profanarlo (1~Mac 2:29--38). El sábado distingue a Israel en cualquier lugar del imperio.
\item \textbf{La sinagoga y la oración.} Sin templo cerca, los judíos de la diáspora se reunían en la \grc{προσευχή} («casa de oración») y en la \grc{συναγωγή}. Inscripciones egipcias del siglo \textsc{iii} a.\,e.\,c. ya mencionan \emph{proseuchai} (Levine, 2005); la inscripción de Teodoto muestra una sinagoga en la misma Jerusalén. Allí se leía la ley cada sábado (Hch 15:21).
\item \textbf{Prosélitos y «temerosos de Dios».} Is 56:6--7 había prometido que los extranjeros que guardan el sábado y se unen a YHWH serían llevados a su santo monte: «mi casa será llamada casa de oración para todos los pueblos» (\heb{כִּי בֵיתִי בֵּית־תְּפִלָּה יִקָּרֵא לְכָל־הָעַמִּים}). En el Segundo Templo hubo prosélitos plenos (\grc{προσήλυτοι}, Mt 23:15; Hch 2:10) y simpatizantes no circuncidados, los «temerosos de Dios» (\grc{φοβούμενοι τὸν θεόν}, Hch 10:2; 13:16), entre los cuales encontrará Pablo su primer auditorio gentil.
\end{itemize}

Zacarías describe la meta de esta vocación con una imagen concreta:
\hebverso{וְהֶחֱזִיקוּ בִּכְנַף אִישׁ יְהוּדִי לֵאמֹר נֵלְכָה עִמָּכֶם כִּי שָׁמַעְנוּ אֱלֹהִים עִמָּכֶם}
\begin{cita}En aquellos días, diez hombres de todas las lenguas de las naciones se asirán, sí, se asirán del manto de un judío, diciendo: «Iremos con vosotros, porque hemos oído que Dios está con vosotros».\fuente{(Zac 8:23)}\end{cita}
El \heb{כָּנָף} (\emph{kanaf}) es el borde del manto, donde se ponían los flecos (\emph{tsitsit}, Nm 15:38): las naciones se acercan al Dios de Israel tomándose de un israelita fiel. Es exactamente la función sacerdotal: mediar el acceso a Dios.

\textbf{La Septuaginta, palabra para las naciones.} La traducción de las Escrituras al griego, la lengua común del Mediterráneo, fue quizá el mayor servicio de Israel a las naciones en esta etapa. Pablo dirá que a los judíos «les fue confiada la palabra de Dios» (Ro 3:2); la Septuaginta es el medio por el que esa palabra llegó a los gentiles. Fue la Biblia de la diáspora, de muchos de los apóstoles y de la iglesia primitiva, y de ella proceden la mayoría de las citas del Antiguo Testamento en el Nuevo.

\textbf{El riesgo: nacionalismo y separación hostil.} La misma etapa muestra cómo la vocación podía quebrarse. La separación necesaria frente a la idolatría (Esd 9--10) podía convertirse en desprecio del otro. Juan Hircano impuso la circuncisión a los idumeos por la fuerza, lo contrario de Is 56. La comunidad de Qumrán se comprometía a «amar a todos los hijos de la luz [\ldots] y odiar a todos los hijos de las tinieblas» (1QS 1:9--10; Vermes, 2011). La balaustrada del templo herodiano, con su advertencia de muerte a los extranjeros, era legítima como protección de la santidad, pero se convirtió en símbolo de división: Pablo dirá que Cristo «derribó el muro intermedio de separación» (\grc{τὸ μεσότοιχον τοῦ φραγμοῦ}, Ef 2:14). La vocación de Israel no era cerrar el acceso, sino ser el lugar donde las naciones encuentran a Dios (Zac 8:23; Is 56:7).

\section{Aporte de la tradición judía}

\subsection{La Gran Asamblea}
La Mishná traza la cadena de la tradición: «Moisés recibió la Torá en el Sinaí y la transmitió a Josué, Josué a los ancianos, los ancianos a los profetas, y los profetas la transmitieron a los hombres de la Gran Asamblea (\heb{אַנְשֵׁי כְנֶסֶת הַגְּדוֹלָה}). Ellos dijeron tres cosas: sed prudentes en el juicio, levantad muchos discípulos y haced un cerco a la Torá» (Avot 1:1). Y añade: «Simón el Justo fue de los últimos de la Gran Asamblea. Él decía: el mundo se sostiene sobre tres cosas: sobre la Torá, sobre el culto y sobre las obras de bondad» (Avot 1:2). La tradición atribuye a esta asamblea, que suele relacionarse con la de Neh 8--10, la fijación de bendiciones y oraciones (Berajot 33a) y el orden de las dieciocho bendiciones de la Amidá, establecidas por «ciento veinte ancianos, entre ellos algunos profetas» (Meguilá 17b). La historicidad de la Gran Asamblea como institución se discute; lo seguro es que la tradición rabínica sitúa en Esdras el comienzo de su propia forma de vida.

\subsection{Esdras, la lectura pública y el cese de la profecía}
El Talmud atribuye a Esdras diez ordenanzas, entre ellas la lectura de la Torá el lunes, el jueves y el sábado por la tarde (Bava Kama 82a), y el paso a la escritura cuadrada aramea (Sanedrín 21b--22a). Lee Neh 8:8 como origen del Targum: «\emph{con claridad}: esto es la traducción; \emph{y ponían el sentido}: la división en versículos» (Meguilá 3a). Son lecturas tradicionales, no datos históricos verificables, pero muestran que el judaísmo rabínico se ve a sí mismo como heredero de Esdras.

La tradición reconoce también el fin de la profecía: «Cuando murieron Hageo, Zacarías y Malaquías, los últimos profetas, el Espíritu Santo se apartó de Israel; con todo, se servían de la \emph{bat qol}» (Tosefta Sotá 13:3; Talmud, Yoma 9b). Seder Olam lo sitúa en tiempos de Alejandro: «hasta aquí los profetas profetizaban por el Espíritu Santo; desde ahora, inclina tu oído y escucha las palabras de los sabios» (Seder Olam 30; Guggenheimer, 1998). Fuentes anteriores coinciden: 1~Mac guarda las piedras del altar profanado «hasta que viniera un profeta» (1~Mac 4:46), habla de una aflicción como no la había «desde el día en que no se vio profeta» (9:27) y nombra a Simón «hasta que se levante un profeta fiel» (14:41); Josefo dice que desde Artajerjes se escribieron historias, pero no con la misma autoridad, «porque no hubo sucesión exacta de profetas» (\emph{Contra Apión} 1.41). Para el Nuevo Testamento, este silencio termina con Juan el Bautista (Lc 3:2; 7:26--28).

\subsection{La reunión de los dispersos en la Amidá}
La décima bendición de la Amidá, que se recita tres veces al día, pide el regreso de todo el pueblo:
\hebverso{תְּקַע בְּשׁוֹפָר גָּדוֹל לְחֵרוּתֵנוּ וְשָׂא נֵס לְקַבֵּץ גָּלֻיּוֹתֵינוּ \foreignlanguage{spanish}{[\ldots]} מְקַבֵּץ נִדְחֵי עַמּוֹ יִשְׂרָאֵל}
\begin{cita}Toca el gran shofar por nuestra libertad, alza estandarte para reunir a nuestros desterrados, y reúnenos juntos de los cuatro extremos de la tierra. Bendito eres tú, YHWH, que reúnes a los dispersos de tu pueblo Israel.\fuente{(Amidá, décima bendición; Sacks, 2009)}\end{cita}
La oración combina Is 27:13 («se tocará el gran shofar») e Is 11:12: «alzará estandarte a las naciones, y reunirá los desterrados de \emph{Israel}, y juntará los esparcidos de \emph{Judá} de los cuatro extremos de la tierra» (\heb{וְאָסַף נִדְחֵי יִשְׂרָאֵל וּנְפֻצוֹת יְהוּדָה יְקַבֵּץ}). En Isaías, «Israel» designa al reino del norte y «Judá» al del sur; la Amidá habla de «tu pueblo Israel» en el sentido del pueblo entero. La redacción actual de la Amidá se fijó después de 70 e.\,c., pero la petición recoge una esperanza que ya estaba viva en el Segundo Templo: el retorno de Babilonia no se consideró la reunión plena de las tribus. Ben Sirá, hacia 180 a.\,e.\,c., pide lo mismo: «reúne a todas las tribus de Jacob» (Eclo 36:10, numeración variable) y espera a Elías «para restablecer las tribus de Jacob» (\grc{καταστῆσαι φυλὰς Ἰακώβ}, Eclo 48:10), uniendo Mal 4:6 [3:24] con Is 49:6.

\subsection{Purim y Janucá}
Las dos fiestas que el judaísmo añadió a las de la Torá nacen en esta etapa, y ambas celebran la preservación del pueblo frente a la amenaza de desaparecer: Purim en la diáspora persa, Janucá en la tierra frente a la helenización forzada. El Talmud explica la ausencia del nombre de Dios en Ester con un juego de palabras sobre Dt 31:18, «yo esconderé (\emph{haster astir}) mi rostro» (Julín 139b): Dios actúa escondido. Sobre Janucá cuenta que los asmoneos hallaron una sola vasija de aceite puro, suficiente para un día, que ardió ocho días (Shabat 21b). Este «milagro del aceite» es tradición talmúdica posterior: 1 y 2~Macabeos no lo mencionan, y explican los ocho días por la dedicación del altar y por la fiesta de Sucot que no se había podido celebrar (2~Mac 10:6). La oración \emph{Al ha-nisim} de Janucá recuerda que «el reino malvado de Grecia se levantó contra tu pueblo Israel para hacerles olvidar tu Torá».

\subsection{Seder Olam y Josefo}
La cronología rabínica tradicional comprime el período persa: Seder Olam le asigna 34 años desde la reconstrucción del templo, y el Talmud reparte los 420 años del Segundo Templo en 34 de Persia, 180 de Grecia, 103 de los asmoneos y 103 de Herodes (Avodá Zará 9a; Guggenheimer, 1998). La historia documentada da a Persia unos dos siglos (539--332). La diferencia se debe a que la tradición identifica a varios reyes persas entre sí; conviene tenerla en cuenta al leer comentarios judíos tradicionales. Para la historia del período, la fuente principal es Flavio Josefo, \emph{Antigüedades judías}, libros 11 a 17, que cubren desde Ciro hasta la muerte de Herodes, con fuentes como la \emph{Carta de Aristeas} y 1~Macabeos.

\section{La promesa de restauración}

\subsection{Un retorno real pero parcial}
Los textos del período reconocen que el retorno no cumplió todas las promesas. La confesión de Neh 9 lo dice con claridad:
\hebverso{הִנֵּה אֲנַחְנוּ הַיּוֹם עֲבָדִים}
\begin{cita}He aquí que hoy somos siervos; en la tierra que diste a nuestros padres para que comieran su fruto y su bien, he aquí que somos siervos en ella. Y su abundante fruto es para los reyes que has puesto sobre nosotros por nuestros pecados [\ldots] y estamos en gran angustia.\fuente{(Neh 9:36--37)}\end{cita}
Faltaban el rey davídico, la reunión de las tribus dispersas por Asiria, la independencia, la gloria visible en el templo y el corazón nuevo prometido por Jeremías y Ezequiel. El libro de Tobías, escrito en el período, lo formula como profecía: los desterrados volverán y reconstruirán el templo, «pero no como el primero, hasta que se cumplan los tiempos» (Tob 14:5).

Dn 9 hace la misma lectura con la profecía de Jeremías. Daniel entiende «por los libros» que los setenta años de desolación de Jerusalén se cumplen (Dn 9:2; Jer 25:11--12; 29:10), ora confesando el pecado de «todo Israel, los de cerca y los de lejos» (Dn 9:7) y recibe esta respuesta:
\hebverso{שָׁבֻעִים שִׁבְעִים נֶחְתַּךְ עַל־עַמְּךָ וְעַל־עִיר קָדְשֶׁךָ}
\begin{cita}Setenta semanas están determinadas sobre tu pueblo y sobre tu santa ciudad, para terminar la prevaricación, poner fin al pecado, expiar la iniquidad, traer la justicia perdurable, sellar la visión y la profecía, y ungir al Santo de los santos.\fuente{(Dn 9:24)}\end{cita}
Los setenta años se convierten en setenta semanas de años. Las interpretaciones de Dn 9:24--27 son muchas: la lectura histórica las cierra en la crisis de Antíoco, con el «ungido» quitado identificado con Onías~III; la lectura mesiánica tradicional cristiana cuenta desde un decreto persa (Esd 7 o Neh 2) hasta el Mesías; otras proyectan la última semana al final de los tiempos (Collins, 1993). Todas coinciden en un punto decisivo para esta serie: el retorno de Ciro no es todavía el final del destierro en sentido pleno. N.~T. Wright (1992) ha sostenido que muchos judíos del Segundo Templo se veían en un «exilio continuado»; la tesis se discute en su alcance, pero textos como Neh 9, Dn 9 y Tob 14 la respaldan en lo esencial.

\subsection{Isaías: la restauración que el retorno no agotó}
Isaías es el profeta que más moldeó la esperanza del período. Sus promesas superan con mucho lo que se vio en 538:

\begin{center}\small
\renewcommand{\arraystretch}{1.3}
\begin{tabular}{>{\RaggedRight\arraybackslash}p{2.4cm}>{\RaggedRight\arraybackslash}p{6.2cm}>{\RaggedRight\arraybackslash}p{5.3cm}}
\toprule
\textbf{Texto} & \textbf{Promesa} & \textbf{En el Segundo Templo y después}\\
\midrule
Is 2:2--4 & Las naciones suben al monte de YHWH para aprender sus caminos & Zac 8:20--22; esperanza no cumplida en el período\\
Is 11:11--12 & Reunión «por segunda vez» de Israel (norte) y Judá (sur) & Amidá, bendición 10; Eclo 36:10\\
Is 44:28--45:1 & Ciro reconstruye Jerusalén y el templo & Esd 1:1--4: cumplimiento histórico\\
Is 49:5--6 & El Siervo levanta las tribus de Jacob y es luz de las naciones & Eclo 48:10; Lc 2:32; Hch 13:47\\
Is 56:6--8 & Extranjeros en el templo; Dios reúne a los dispersos de Israel «y a otros» & Prosélitos; Mc 11:17\\
Is 60--62 & Las naciones traen a los hijos de Sion; Israel, «sacerdotes de YHWH» (61:6) & Lc 4:18--21; Ro 11:12, 15\\
Is 66:18--23 & Mensajeros a las naciones; de ellas, sacerdotes y levitas & Ro 15:16; misión de Pablo\\
\bottomrule
\end{tabular}\end{center}

Dos textos merecen atención especial. En Is 49:5--6, el Siervo recibe primero la misión de «levantar las tribus de Jacob y restaurar a los preservados de Israel», y luego una misión mayor: «poco es para mí que solo seas mi siervo para [\ldots] restaurar el remanente de Israel; también te di por luz de las naciones (\heb{וּנְתַתִּיךָ לְאוֹר גּוֹיִם}), para que seas mi salvación hasta lo último de la tierra». La restauración de las tribus y la luz para las naciones forman una sola obra; Simeón citará este texto al tomar al niño Jesús (Lc 2:32). En Is 56:8, el Dios «que reúne a los dispersos de Israel» (\heb{מְקַבֵּץ נִדְחֵי יִשְׂרָאֵל}, la misma fórmula de la Amidá) promete: «aún juntaré sobre él a otros, además de sus reunidos». Israel reunido y las naciones añadidas: es la estructura que Pablo desarrollará con la imagen del olivo (Ro 11:17--24).

Is 60--62 describe a Sion reconstruida: «andarán las naciones a tu luz» (Is 60:3), «edificarán las ruinas antiguas» (61:4), «vosotros seréis llamados sacerdotes de YHWH» (61:6), y sobre los muros de Jerusalén Dios pone centinelas que no callan «hasta que establezca a Jerusalén y la ponga por alabanza en la tierra» (62:6--7). Nehemías construyó los muros; la alabanza universal quedó pendiente. Is 66:18--23 cierra el libro con mensajeros enviados a Tarsis, Pul (o Put), Lud, Tubal y Javán (Grecia), que traen a «vuestros hermanos de todas las naciones» como ofrenda, y con sacerdotes y levitas tomados de ellas (66:21). Pablo entenderá su propio ministerio como el de un sacerdote que ofrece «la ofrenda de los gentiles» (Ro 15:16).

\subsection{Jeremías: setenta años y algo más}
Jeremías dio la fecha del retorno y, a la vez, una promesa que lo excede:
\hebverso{ושבתי את שבותכם וקבצתי אתכם מכל הגוים}
\begin{cita}Cuando se cumplan en Babilonia los setenta años, yo os visitaré [\ldots]. Me buscaréis y me hallaréis, porque me buscaréis de todo vuestro corazón. Y seré hallado por vosotros, dice YHWH, y haré volver vuestra cautividad, y os reuniré de todas las naciones y de todos los lugares adonde os arrojé.\fuente{(Jer 29:10, 13--14)}\end{cita}
El edicto de Ciro cumplió la primera parte («setenta años»); la segunda («de todas las naciones», «de todo vuestro corazón») abrió el horizonte que Dn 9 reinterpreta. Jer 33:14--18 promete cumplir «la buena palabra que he hablado a la casa de Israel y a la casa de Judá», con un Renuevo de David y la promesa de que «no faltará a David varón que se siente sobre el trono de la casa de Israel» (\heb{לֹא־יִכָּרֵת לְדָוִד אִישׁ יֹשֵׁב עַל־כִּסֵּא בֵית־יִשְׂרָאֵל}, 33:17). En el Segundo Templo no se sentó ningún davídico en el trono: Zorobabel fue gobernador, los asmoneos eran levitas y Herodes idumeo. La promesa quedó abierta y se concentró en la espera del Hijo de David. Jeremías había anunciado también que «en aquellos tiempos irá la casa de Judá a la casa de Israel, y vendrán juntamente de la tierra del norte» (Jer 3:18; cf. 50:4--5): la reunión de las dos casas tampoco ocurrió con Ciro.

\subsection{La esperanza mesiánica al final del período}
Tras la conquista de Pompeyo, los \emph{Salmos de Salomón}, escritos en círculos piadosos de Jerusalén, piden: «Mira, Señor, y levántales su rey, el hijo de David [\ldots] para que purifique a Jerusalén de las naciones», que «reunirá un pueblo santo» y «distribuirá las tribus sobre la tierra» (Sal Salomón 17:21--22, 26, 28; R.~B. Wright, 1985). La esperanza davídica, la santidad del pueblo y las doce tribus aparecen unidas. En Qumrán, la \emph{Regla de la comunidad} espera «la venida de un profeta y de los mesías de Aarón y de Israel» (1QS 9:11); el \emph{Florilegio} (4Q174) aplica 2~S 7 y Am 9:11 («la tienda de David caída») al Mesías; el texto 11QMelquisedec (11Q13) une Is 61:1--2, el jubileo de Lv 25 y la cronología de Dn 9 para anunciar una liberación final; y 4Q521 describe al ungido a quien «obedecerán los cielos y la tierra», que libera a los cautivos, da vista a los ciegos y resucita muertos (Vermes, 2011; VanderKam, 2010). Lo que Jesús responde a los discípulos de Juan (Mt 11:5) se lee sobre este trasfondo.

\subsection{Proyección en el Nuevo Testamento y en Romanos 11}
Lucas presenta a quienes esperaban la restauración en el templo de Herodes:
\begin{cita}Este hombre, justo y piadoso, esperaba la consolación de Israel (\upshape\grc{προσδεχόμενος παράκλησιν τοῦ Ἰσραήλ}\itshape), y el Espíritu Santo estaba sobre él.\fuente{(Lc 2:25)}\end{cita}
«Consolación» remite a Is 40:1 («consolad, consolad a mi pueblo», LXX \grc{παρακαλεῖτε}). Ana, de la tribu de Aser, hablaba del niño «a todos los que esperaban la redención de Jerusalén» (\grc{λύτρωσιν Ἰερουσαλήμ}, Lc 2:38). Al final del evangelio, los discípulos de Emaús dicen: «esperábamos que él era el que había de redimir a Israel» (Lc 24:21), y en Hch 1:6 preguntan: «Señor, ¿restaurarás el reino a Israel en este tiempo?» (\grc{ἀποκαθιστάνεις τὴν βασιλείαν τῷ Ἰσραήλ}). Jesús no niega la restauración; corrige el calendario («no os toca saber los tiempos») y abre la misión hasta lo último de la tierra (Hch 1:7--8), en el orden de Is 49:6. Pablo, ante Agripa, habla de la promesa «que nuestras doce tribus (\grc{τὸ δωδεκάφυλον ἡμῶν}), sirviendo constantemente a Dios de día y de noche, esperan alcanzar» (Hch 26:7): para él, las doce tribus siguen existiendo como sujeto de la esperanza.

Pablo resume el sentido de esta etapa en una frase: «cuando vino el cumplimiento del tiempo (\grc{τὸ πλήρωμα τοῦ χρόνου}), Dios envió a su Hijo, nacido de mujer, nacido bajo la ley» (Gá 4:4). Los siglos del Segundo Templo son el tiempo que se «llena» antes del Mesías. La misma palabra, \grc{πλήρωμα}, reaparecerá en Ro 11:12 («la plenitud de ellos») y en Ro 11:25 («la plenitud de los gentiles»). Y cuando Pablo anuncie que «todo Israel será salvo», lo fundará en Is 59:20--21 y 27:9, con eco del nuevo pacto de Jer 31:33--34 (Ro 11:26--27), las promesas que el retorno de Ciro dejó abiertas. ONE FOR ISRAEL (2025) lee Ro 11:26 como un futuro volverse del pueblo judío a Yeshúa como Mesías, con los gentiles creyentes injertados en un Israel «ampliado, no reemplazado»; el Estudio~8 examinará esa y las demás lecturas.

\section{Conclusión y puente al Estudio 7}

El período del Segundo Templo deja cuatro resultados para la serie.

\begin{enumerate}
\item \textbf{El retorno fue real.} Dios cumplió la palabra de Jeremías por medio de Ciro: Judá y Benjamín, con sacerdotes, levitas y miembros de otras tribus, volvieron, reconstruyeron el templo y ofrecieron por las doce tribus (Esd 6:17).
\item \textbf{El retorno fue parcial.} Sin rey davídico, sin las diez tribus dispersas por Asiria como cuerpo, bajo imperios extranjeros y con la gloria prometida aún por venir, los mismos libros del período lo reconocen (Neh 9:36--37; Dn 9; Tob 14:5). Los profetas siguieron hablando de «la casa de Judá y la casa de Israel» (Zac 8:13; Jer 33:14).
\item \textbf{Israel aprendió a ser pueblo santo sin trono.} Torá, sábado, sinagoga, oración y Septuaginta sostuvieron la identidad del pueblo y lo abrieron a prosélitos y temerosos de Dios, aunque también se manifestó el riesgo de la separación hostil.
\item \textbf{La esperanza quedó concentrada.} Malaquías anuncia al Señor que viene a su templo y a Elías que restaura; Isaías y Jeremías, la reunión de las tribus y la luz para las naciones; Qumrán y los Salmos de Salomón, al Hijo de David.
\end{enumerate}

Al final del período, en el templo que Herodes estaba ampliando, un sacerdote llamado Zacarías recibe el anuncio de un hijo que irá delante del Señor «con el espíritu y el poder de Elías, para hacer volver los corazones de los padres a los hijos» (Lc 1:17; Mal 4:6 [3:24]). Con esa cita de Malaquías, el silencio termina. El Estudio~7, \emph{de Cristo a la destrucción del templo}, seguirá al Mesías en su contexto judío, al remanente que creyó y a la caída del templo en 70 e.\,c.

\section*{Preguntas para la reflexión}\markboth{Preguntas para la reflexión}{}
\begin{enumerate}
\item La comunidad que volvió era sobre todo de Judá y Benjamín, pero ofreció doce machos cabríos «por todo Israel» (Esd 6:17). ¿Qué enseña este gesto sobre la relación entre un remanente y el pueblo entero (Ro 11:5, 16)?
\item Nehemías ora «hoy somos siervos» en plena restauración (Neh 9:36). ¿Cómo distinguir entre un cumplimiento real y un cumplimiento pleno de las promesas de Dios?
\item Zac 8:23 imagina a las naciones tomadas del manto de un judío. ¿Qué actitudes de Israel favorecían esa escena y cuáles la impedían en el Segundo Templo? ¿Qué dice esto a los creyentes de las naciones (Ro 11:18)?
\item ¿Qué cambia en la lectura de Hch 1:6 cuando se recuerda que la Amidá pedía cada día la reunión de los dispersos de Israel?
\end{enumerate}

\section*{Glosario}\markboth{Glosario}{}
\begin{description}
\item[Amaná] (\heb{אֲמָנָה}). Pacto firmado o compromiso escrito; el de Neh 9:38 [10:1].
\item[Anshei Knéset ha-Guedolá] (\heb{אַנְשֵׁי כְנֶסֶת הַגְּדוֹלָה}). «Hombres de la Gran Asamblea». Cuerpo que la tradición rabínica sitúa entre los profetas y los sabios (Avot 1:1).
\item[Bat qol] (\heb{בַּת קוֹל}). «Hija de la voz». Voz del cielo que, según la tradición, sustituye a la profecía tras Malaquías.
\item[Diáspora] (\grc{διασπορά}). Dispersión del pueblo judío fuera de la tierra (Dt 28:25 LXX; Jn 7:35; Stg 1:1).
\item[Enkainia / Janucá] (\grc{τὰ ἐγκαίνια}; \heb{חֲנֻכָּה}). «Dedicación». Fiesta de la purificación del templo en 164 a.\,e.\,c. (1~Mac 4:59; Jn 10:22).
\item[Mashíaj] (\heb{מָשִׁיחַ}). Ungido. Título dado a Ciro en Is 45:1; en el período, esperanza del ungido davídico.
\item[Mityahadim] (\heb{מִתְיַהֲדִים}). «Se hacían judíos» (Est 8:17). Primera mención bíblica de gentiles que se unen al pueblo judío.
\item[Pleroma] (\grc{πλήρωμα}). Plenitud, cumplimiento. «El cumplimiento del tiempo» (Gá 4:4); clave en Ro 11:12, 25.
\item[Prosélito] (\grc{προσήλυτος}). Extranjero incorporado plenamente a Israel; traduce el hebreo \heb{גֵּר} (\emph{guer}) en la LXX.
\item[Shiqquts meshomem] (\heb{שִׁקּוּץ מְשׁוֹמֵם}). «Abominación desoladora» (Dn 11:31); en griego \grc{βδέλυγμα ἐρημώσεως} (1~Mac 1:54; Mt 24:15).
\item[Sofer] (\heb{סוֹפֵר}). Escriba, experto en la Torá. Título de Esdras (Esd 7:6).
\item[Yehud / Yehudí] (\heb{יְהוּד}, \heb{יְהוּדִי}). Provincia persa de Judá; «judío», de Judá y luego de todo el pueblo (Est 2:5).
\end{description}

\section*{Referencias}\markboth{Referencias}{}
\refapa{Bickerman, E.~J. (1988). \emph{The Jews in the Greek age}. Harvard University Press.}
\refapa{\emph{Biblia Hebraica Stuttgartensia} (5.ª ed.). (1997). Deutsche Bibelgesellschaft.}
\refapa{Charlesworth, J.~H. (Ed.). (1985). \emph{The Old Testament pseudepigrapha} (Vol. 2). Doubleday.}
\refapa{Collins, J.~J. (1993). \emph{Daniel} (Hermeneia). Fortress Press.}
\refapa{Cotton, H.~M., \& Wörrle, M. (2007). Seleukos IV to Heliodoros: A new dossier of royal correspondence from Israel. \emph{Zeitschrift für Papyrologie und Epigraphik, 159}, 191--205.}
\refapa{Crown, A.~D. (Ed.). (1989). \emph{The Samaritans}. Mohr Siebeck.}
\refapa{Danby, H. (Trad.). (1933). \emph{The Mishnah}. Oxford University Press.}
\refapa{Departamento de Educación, Congregación Iglesia Israelita del Nuevo Pacto. (s.\,f.). \emph{Restauración de Israel} (Clase 7) [Presentación].}
\refapa{Epstein, I. (Ed.). (1935--1948). \emph{The Babylonian Talmud}. Soncino Press.}
\refapa{Finkel, I. (Ed.). (2013). \emph{The Cyrus Cylinder: The king of Persia's proclamation from ancient Babylon}. I.\,B. Tauris.}
\refapa{Finkelstein, I. (2008). Jerusalem in the Persian (and early Hellenistic) period and the wall of Nehemiah. \emph{Journal for the Study of the Old Testament, 32}(4), 501--520.}
\refapa{Goldstein, J.~A. (1976). \emph{I Maccabees} (Anchor Bible 41). Doubleday.}
\refapa{Guggenheimer, H.~W. (1998). \emph{Seder Olam: The rabbinic view of biblical chronology}. Jason Aronson.}
\refapa{\emph{Israel, el pueblo de Dios} (Octava lección). (s.\,f.). [Material de estudio].}
\refapa{Jobes, K.~H., \& Silva, M. (2015). \emph{Invitation to the Septuagint} (2.ª ed.). Baker Academic.}
\refapa{Josefo, F. (1987). \emph{The works of Josephus} (W. Whiston, Trad.). Hendrickson.}
\refapa{Kloppenborg, J.~S. (2000). Dating Theodotos (CIJ II 1404). \emph{Journal of Jewish Studies, 51}(2), 243--280.}
\refapa{Kuhrt, A. (1983). The Cyrus Cylinder and Achaemenid imperial policy. \emph{Journal for the Study of the Old Testament, 8}(25), 83--97.}
\refapa{Levine, L.~I. (2005). \emph{The ancient synagogue: The first thousand years} (2.ª ed.). Yale University Press.}
\refapa{Lipschits, O., \& Vanderhooft, D.~S. (2011). \emph{The Yehud stamp impressions: A corpus of inscribed impressions from the Persian and Hellenistic periods in Judah}. Eisenbrauns.}
\refapa{Magen, Y. (2008). \emph{Mount Gerizim excavations: Vol. 2. A temple city}. Israel Antiquities Authority.}
\refapa{Mazar, B. (1975). \emph{The mountain of the Lord}. Doubleday.}
\refapa{Meshorer, Y. (2001). \emph{A treasury of Jewish coins: From the Persian period to Bar Kokhba}. Amphora Books; Yad Ben-Zvi Press.}
\refapa{Neusner, J. (Trad.). (2002). \emph{The Tosefta}. Hendrickson.}
\refapa{\emph{Novum Testamentum Graece} (28.ª ed.). (2012). Deutsche Bibelgesellschaft.}
\refapa{ONE FOR ISRAEL Ministry. (2025, 3 de julio). \emph{What does «all Israel will be saved» mean in Romans 11?} \url{https://www.oneforisrael.org/what-does-all-israel-will-be-saved-mean-in-romans-11/}}
\refapa{Pearce, L.~E., \& Wunsch, C. (2014). \emph{Documents of Judean exiles and West Semitic exiles in Babylonia in the collection of David Sofer} (CUSAS 28). CDL Press.}
\refapa{Porten, B., \& Yardeni, A. (1986). \emph{Textbook of Aramaic documents from ancient Egypt: Vol. 1. Letters}. Hebrew University.}
\refapa{Rahlfs, A., \& Hanhart, R. (Eds.). (2006). \emph{Septuaginta} (ed. rev.). Deutsche Bibelgesellschaft.}
\refapa{Sacks, J. (Trad.). (2009). \emph{The Koren Siddur}. Koren Publishers.}
\refapa{Schürer, E. (1973--1987). \emph{The history of the Jewish people in the age of Jesus Christ} (G. Vermes, F. Millar y M. Goodman, Eds. rev.; 3 vols.). T\&T Clark.}
\refapa{Ulrich, E., \& Flint, P.~W. (2010). \emph{Qumran Cave 1.II: The Isaiah scrolls} (Discoveries in the Judaean Desert 32). Clarendon Press.}
\refapa{VanderKam, J.~C. (2010). \emph{The Dead Sea scrolls today} (2.ª ed.). Eerdmans.}
\refapa{Vermes, G. (2011). \emph{The complete Dead Sea scrolls in English} (ed. rev.). Penguin.}
\refapa{Williamson, H.~G.~M. (1985). \emph{Ezra, Nehemiah} (Word Biblical Commentary 16). Word Books.}
\refapa{Wright, N.~T. (1992). \emph{The New Testament and the people of God}. Fortress Press.}
\refapa{Wright, R.~B. (1985). Psalms of Solomon. En J.~H. Charlesworth (Ed.), \emph{The Old Testament pseudepigrapha} (Vol. 2). Doubleday.}
\end{document}
